\appendix
\section{Validação da estratégia de busca}
\label{ap:validacao-busca}

Este apêndice documenta o processo de validação aplicado à estratégia de busca antes da execução das consultas finais nas bases.
A validação combina dois procedimentos: auditoria das strings contra revisões sistemáticas comparáveis e verificação de recall por meio de um conjunto de artigos-semente operando como Quasi-Gold Standard (QGS).
A última subseção apresenta o reporte completo do PRISMA-S~\cite{rethlefsen2021prismas}.

\subsection{Auditoria pré-execução}
\label{ap:auditoria}

A auditoria pré-execução consultou seis revisões sistemáticas e mapeamentos publicados entre 2024 e 2026 sobre agentes de codificação baseados em \textit{Large Language Models} (LLM), agentes de software autônomos e técnicas de memória em agentes de IA.
Para cada revisão foram extraídos: as strings completas (quando disponíveis), o conjunto de termos do bloco de agentes e do bloco de memória, e os sistemas representativos catalogados.
Em paralelo, foi compilada uma lista de 16 sistemas conhecidos da literatura corrente como casos-teste de recall, e dez formulações alternativas de termos foram testadas individualmente em consultas exploratórias no Scopus e no arXiv.
A \autoref{tab:auditoria-slrs} resume as revisões consultadas e a contribuição de cada uma para o refinamento dos blocos.

\begin{table}[htbp]
\centering
\caption{Revisões sistemáticas auditadas para calibração das strings de busca.}
\label{tab:auditoria-slrs}
\small
\begin{tabular}{p{3.2cm}p{1cm}p{2.2cm}p{1cm}p{4.5cm}}
\toprule
\textbf{Revisão} & \textbf{Ano} & \textbf{Venue} & \textbf{N} & \textbf{Contribuição ao refinamento} \\
\midrule
Hou et al.~\cite{hou2024llmse} & 2024 & TOSEM & 395 & Termos de bloco \textit{coding}/\textit{software engineering}; cobertura de venues. \\
Liu et al.~\cite{liu2024llmagents} & 2024 & TOSEM & 117 & Vocabulário de agentes de software autônomos. \\
Zhang et al.~\cite{zhang2024unifying} & 2024 & SCIS & 52 & Termos de planejamento e memória em agentes LLM. \\
Jin et al.~\cite{jin2024llmagents} & 2024 & arXiv & 90 & Variantes de \textit{LLM-based agent}; sinônimos de bloco C1. \\
Wang et al.~\cite{wang2024llmagents} & 2024 & ASE & 106 & Confirmação do termo \textit{language agent}; vocabulário de papéis. \\
Du et al.~\cite{du2026memory} & 2026 & arXiv & --- & Taxonomia de memória em agentes de IA; vocabulário de bloco C2. \\
\bottomrule
\end{tabular}
\end{table}

A auditoria identificou três lacunas no rascunho inicial das strings: ausência de variantes sem sufixo \textit{-ing} (\textit{code assistant}), ausência do termo emergente \textit{agentic coding} (já adotado como categoria de venue em 2025) e ausência da forma \textit{SWE agent} como nome próprio de classe.
As lacunas correspondentes no bloco de memória abrangiam \textit{persistent memory}, \textit{persistent context} e \textit{cross-task}, recorrentes nos papers-semente e nas revisões auditadas.

\subsection{Recuperação de artigos-semente (Quasi-Gold Standard)}
\label{ap:qgs}

A validação por Quasi-Gold Standard consiste em verificar se uma estratégia de busca recupera um conjunto pré-selecionado de artigos representativos do escopo, identificados independentemente da estratégia sendo validada~\cite{mourao2020hybrid}.
Foram definidos sete artigos-semente, escolhidos por cobrirem diferentes arquiteturas de persistência (banco de experiência, memória episódica, memória orientada a commits, contexto inspirado em \textit{git}) e por incluírem ao menos um resultado negativo.
A estratégia passou por dois ciclos de refinamento documentados no
protocolo versionado, elevando a recuperação de quatro para seis dos
sete artigos-semente; o caso restante (CodeMEM) foi resgatado pelo rastreamento de citações, conforme descrito na seção de método.
A \autoref{tab:qgs-recuperacao} consolida o status de recuperação por base.

\begin{table}[htbp]
\centering
\caption{Recuperação dos artigos-semente por base e método. \textit{Cit.} indica resgate via rastreamento de citações.}
\label{tab:qgs-recuperacao}
\small
\begin{tabular}{p{3.5cm}p{1cm}p{1.5cm}p{1.2cm}p{1.5cm}p{2cm}}
\toprule
\textbf{Artigo-semente} & \textbf{Ano} & \textbf{Scopus} & \textbf{arXiv} & \textbf{S2 ID} & \textbf{Recuperado} \\
\midrule
Du (survey de memória)~\cite{du2026memory} & 2026 & não & sim & sim & Sim (string) \\
Joshi et al.\ (SWE-Bench-CL)~\cite{joshi2025swebenchcl} & 2025 & não & sim & sim & Sim (string) \\
Wang et al.\ (CodeMEM)~\cite{wang2026codemem} & 2026 & não & não & sim & Sim (Cit.) \\
Deng et al.\ (MemCoder)~\cite{deng2026memcoder} & 2026 & não & sim & sim & Sim (string) \\
Wu et al.\ (GCC)~\cite{wu2025gcc} & 2025 & não & sim & sim & Sim (string) \\
Wang et al.\ (Repository Memory)~\cite{wang2026improving} & 2026 & não & sim & sim & Sim (string) \\
Lindenbauer et al.\ (CTIM-Rover)~\cite{lindenbauer2025ctimrover} & 2025 & não & sim & sim & Sim (string) \\
\bottomrule
\end{tabular}
\end{table}

A ausência total dos artigos-semente no Scopus reflete uma característica do campo: a maioria dos sistemas analisados é composta de \textit{preprints} de 2025--2026 ainda fora do ciclo de indexação de venues tradicionais.
Esse diagnóstico foi confirmado por consultas auxiliares no Scopus (busca por nomes próprios de sistemas), que retornaram zero resultados, e justifica a centralidade do rastreamento de citações como terceiro método PRISMA na estratégia adotada.

\subsection{Conformidade PRISMA-S}
\label{ap:prisma-s}

A \autoref{tab:prisma-s} apresenta o reporte da estratégia de busca segundo os 16 itens da extensão PRISMA-S~\cite{rethlefsen2021prismas} ao PRISMA 2020~\cite{page2021prisma}.
Os itens cobrem nomeação das fontes, plataformas, limites, métodos de busca complementares, gerenciamento de referências e atualizações.

\begin{longtable}{p{0.5cm}p{4cm}p{9cm}}
\caption{Reporte PRISMA-S (16 itens) da estratégia de busca.}
\label{tab:prisma-s} \\
\toprule
\textbf{\#} & \textbf{Item} & \textbf{Descrição} \\
\midrule
\endfirsthead
\toprule
\textbf{\#} & \textbf{Item} & \textbf{Descrição} \\
\midrule
\endhead
\bottomrule
\endfoot
1 & Nomes das bases & Scopus e arXiv. \\
2 & Plataforma/interface & Scopus via API Elsevier (\texttt{scopus-mcp}); arXiv via API REST com paginação completa (\texttt{fetch\_arxiv.py}). \\
3 & Registros de estudos & Nenhum registro de estudos (por exemplo, PROSPERO) consultado. \\
4 & Recursos online & Nenhum site organizacional consultado separadamente. \\
5 & Rastreamento de citações & Forward e backward via Semantic Scholar a partir de sete artigos-semente; ResearchRabbit como alternativa manual. \\
6 & Contato com autores & Nenhum autor contactado. \\
7 & Outros métodos & Nenhum além dos descritos nos itens 1, 2 e 5. \\
8 & Estratégias completas & Strings exatas por base documentadas no protocolo de busca, com versionamento em \textit{git}. \\
9 & Limites e restrições & Scopus: \texttt{PUBYEAR > 2022}; arXiv: categorias \texttt{cs.SE}, \texttt{cs.AI}, \texttt{cs.CL} e filtro pós-busca \texttt{--year-from 2023}; idiomas inglês e português. \\
10 & Desenvolvimento da estratégia & Mapeamento de três blocos conceituais, construção das strings com OR/AND, validação por sete artigos-semente e dois ciclos de refinamento iterativo documentados no protocolo versionado. \\
11 & Revisão por pares & Estratégia revisada pelo orientador antes da execução final. \\
12 & Gerenciamento de registros & Zotero com Better BibTeX para gerenciamento; arquivos \texttt{.bib} versionados em \textit{git}. \\
13 & Total de registros & 1.007 dos bancos (Scopus: 39; arXiv: 968) e 186 únicos por rastreamento de citações; total de 1.193 antes da deduplicação. \\
14 & Desduplicação & Elicit Pro deduplica automaticamente na importação; remoção manual subsequente totalizando 83 duplicatas. \\
15 & Ferramentas de automação & Scopus via MCP, arXiv via script Python, triagem assistida por Elicit Pro, rastreamento via Semantic Scholar MCP. \\
16 & Atualizações & Busca a ser re-executada antes da submissão se decorridos mais de três meses desde a execução inicial. \\
\end{longtable}
